\subsection{Erstellung des Aufgabenkatalogs für den Usability-Test}
\label{ch:Validierung:sec:Material:Usability_Test}

Der Usability-Test wurde auf Grundlage des Fragebogens gestaltet, um alle abgefragten Szenarien abzudecken und Unterschiede zwischen den Systemen für die Bewertung hervorzuheben.

Der daraus resultierende Aufgabenkatalog ist in \autoref{Add:Usability_Test:Aufgabenkatalog} aufgelistet. Er umfasst den täglich benötigten allgemeinen Status der Tests, die Abfrage verschiedener Varianten von absoluten Werten, dessen Ergebnisse jeweils unterschiedlich abzulesen sind, sowie den Wechsel zwischen den beiden Dash\-boards. Ein zusammenfassender Überblick über die Aufgabentypen ist in \autoref{tab:usability_test:Überblick_Aufgabentypen} aufgelistet. Eine detaillierte Auflistung über die Gründe der Aufgabenauswahl und der verwendeten Reihenfolge ist in \autoref{Add:Usability_Test:Begündung_Filterung} aufgeführt.

Nachfolgend wird mit der Kurzform 'Aufgabe 1', 'Aufgabe 2' etc. auf den Aufgabenkatalog des Usability-Tests referenziert.

\begin{table}[]
    \begin{tabular}{|c|c|cc|}
    \hline
    \multirow{2}{*}{Nr} & \multirow{2}{*}{Stichwort} & \multicolumn{2}{c|}{Aufgabentyp}                                                                                                                                                                                           \\ \cline{3-4} 
                        &                            & \multicolumn{1}{c|}{altes System}                                                                                    & neues System                                                                                        \\ \hline
    \multirow{2}{*}{1}  & \multirow{2}{*}{Green Day} & \multicolumn{1}{c|}{\begin{tabular}[c]{@{}c@{}}sind rote Felder mit 'ERR'\\ oder Zahlen != 0 zu sehen?\end{tabular}} & Sind rote Zahlen zu sehen?                                                                          \\ \cline{3-4} 
                        &                            & \multicolumn{2}{c|}{Zusammenfassung von 2 Systemen}                                                                                                                                                                        \\ \hline
    \multirow{2}{*}{2}  & \multirow{2}{*}{Build}     & \multicolumn{1}{c|}{\begin{tabular}[c]{@{}c@{}}Rechenaufgabe\\ (ohne Farbgestaltung)\end{tabular}}                   & \begin{tabular}[c]{@{}c@{}}Ableseaufgabe\\ (rote Zahl von 2 Systemen\\ zusammenzählen)\end{tabular} \\ \cline{3-4} 
                        &                            & \multicolumn{2}{c|}{Zusammenfassung von 2 Systemen}                                                                                                                                                                        \\ \hline
    \multirow{2}{*}{3}  & \multirow{2}{*}{Release}   & \multicolumn{1}{c|}{\begin{tabular}[c]{@{}c@{}}Rechenaufgabe\\ (2 neue Oberflächen)\end{tabular}}                    & \begin{tabular}[c]{@{}c@{}}Ableseaufgabe\\ (2. Dashbard)\end{tabular}                               \\ \cline{3-4} 
                        &                            & \multicolumn{2}{l|}{\begin{tabular}[c]{@{}l@{}}- das System / Oberfläche wird gewechselt\\ - je eine Aufgabe pro System\end{tabular}}                                                                                      \\ \hline
    \multirow{2}{*}{4}  & \multirow{2}{*}{Modi}      & \multicolumn{1}{c|}{Zählaufgabe}                                                                                     & Ableseaufgabe                                                                                       \\ \cline{3-4} 
                        &                            & \multicolumn{2}{l|}{\begin{tabular}[c]{@{}l@{}}- das System / Oberfläche wird zurück gewechselt\\ - je eine Aufgabe pro System\end{tabular}}                                                                               \\ \hline
    \end{tabular}
    \caption{Usability-Test: Überblick über die Aufgabentypen (Zusammenfassung aus \autoref{Add:Usability_Test:Aufgabenkatalog} und \autoref{Add:Usability_Test:Begündung_Filterung})}
    \label{tab:usability_test:Überblick_Aufgabentypen}
\end{table}